\hypertarget{the-filesystem}{%
\chapter{6. The filesystem}\label{the-filesystem}}

Part I built the machinery: the kernel facilities woven into a
confinement, the processes that hold it, the bus, the trusted base, and
the surface that drives it. Part II turns to the resource classes that
machinery confines, a mechanism at a time, and it starts with the
filesystem, because the filesystem is where the project's motivation
lives. The thing a kennel must not see is the user's real home
directory, with its keys, its tokens, and its shell history, and the
account of how a kennel does not see it is the account of the
filesystem.

\hypertarget{the-view-and-the-access}{%
\section{6.1 The view and the access}\label{the-view-and-the-access}}

Filesystem confinement is built from both limbs at once, and the design
is explicit that neither alone is enough. The construction limb is the
view: a mount namespace and bind mounts assemble what the kennel sees in
its home and runtime directories out of only the granted paths, so the
real directory contents are not there to be reached
(\texttt{construction-by-absence}). The interposition limb is Landlock,
a kernel-enforced ACL keyed to the process that grants read, write, and
execute on the bound paths and denies the rest
(\texttt{interpose-as-transaction}).

Each closes a gap the other leaves open. A view without an ACL can be
read around: a kennel could follow a symlink out of a bound directory
into the real path behind it, and with no ACL in the kernel nothing
would stop the dereference. An ACL without a view leaks by enumeration:
a \texttt{readdir} on the home directory would return every entry,
including the ones the kennel cannot open, which both tells the workload
what exists and confuses the user about what was granted. With both, a
\texttt{readdir} returns only the granted entries because the others are
absent from the view, and a symlink that reaches for a real path is
resolved by Landlock to that path and denied there. The symlink class is
the kernel's to handle: Landlock resolves a link to its target and
applies the ruleset to the resolved path, so a link out of a granted
directory gains nothing. The symmetric bind-mount escape, remounting a
forbidden path into an allowed one, is closed by the kennel holding no
\texttt{CAP\_SYS\_ADMIN} to mount with, by \texttt{MS\_NODEV},
\texttt{MS\_NOSUID}, and \texttt{MS\_NOEXEC} on the view's mount points,
and by the namespace being \texttt{MS\_SLAVE} so a mount inside it
cannot propagate back to the host.

\hypertarget{the-constructed-home}{%
\section{6.2 The constructed home}\label{the-constructed-home}}

The home directory is the heart of the construction, and the move is to
build a new one. The view begins as a fresh tmpfs that becomes the
kennel's root, and into it the spawn binds the policy's granted paths
from the real home, read-only unless the grant is writable. The grant
names the paths and which of them may be written:

\begin{kennelcode}
[fs]
read  = ["~/projects/acme/**", "~/.config/acme/config.toml"]
write = ["~/projects/acme/build/**"]
\end{kennelcode}

Each read path binds into the fresh root read-only and each write path
binds writable, and the home those binds land in otherwise holds
nothing:

\begin{kennelcode}
~/                              fresh tmpfs, writable
  projects/acme/                bound read-only
  projects/acme/build/          bound writable, persists to the host
  .config/acme/config.toml      bound read-only
\end{kennelcode}

The system surfaces around them are built rather than shared:
\texttt{/etc}'s identity files are written scrubbed of host specifics,
naming an in-kennel user, so a lookup of the running identity returns
the masked name and the constructed home rather than the operator's real
login.

A policy names home as \texttt{\textasciitilde{}}, and only as
\texttt{\textasciitilde{}}. The same
\texttt{\textasciitilde{}/projects/acme} is the operator's real
\texttt{/home/johndoe/projects/acme} on the host, where the bind source
lives, and the kennel's own \texttt{/home/kennel/projects/acme} inside
the view, where it lands, under the home of the masked \texttt{kennel}
user. The two absolute paths differ, and \texttt{\textasciitilde{}} is
the one token that resolves to the right home on each side of the bind.
A literal host path in a grant would pin one operator's login into a
document meant to be portable, and would name the host's home where the
view's was meant, so a grant never writes one. Home is
\texttt{\textasciitilde{}}, on both sides of the boundary, always.

The spawn then \texttt{pivot\_root}s into the new root, and the Landlock
ruleset is applied afterward so its rules key on the view's own inodes.

The result is the protection the project was built for. Inside the
kennel, a listing of the home shows only the granted entries; a path to
the user's SSH directory returns no such file, because the directory was
never built into the view; and a constructed path that tries to reach
the real one is denied by Landlock as well, so the absence is backed by
a denial. The kennel cannot enumerate, cannot discover, and cannot
stumble into the credentials and state in the user's real home, because
they are not in the world it was given
(\texttt{construction-by-absence}). The protection is absence rather
than redaction: the credential paths are not scrubbed in place, they are
simply not bound, which is what closes the casual-reconnaissance case
(T1.1) without any in-view filtering at all.

\hypertarget{writable-but-ephemeral}{%
\section{6.3 Writable, but ephemeral}\label{writable-but-ephemeral}}

The constructed home carries a Landlock write grant, because an
application that cannot write to its own home breaks in surprising ways,
and least astonishment is worth a grant here. What makes that safe is
which surface is writable: it is the fresh tmpfs, not the host. Anything
the workload creates directly in its home lands on that ephemeral tmpfs
and is gone at teardown; only the paths bound writable from the real
home survive, and a path bound read-only stays read-only at the VFS
layer regardless of the home's own write grant. So the writable home is
real enough for tools that expect it and empty of consequence by
default: a workload that scatters files across its home persists nothing
it was not explicitly granted a writable bind for, and the persistence
it does get is exactly the set of paths the policy named.

\hypertarget{denials-and-the-pseudo-filesystems}{%
\section{6.4 Denials and the
pseudo-filesystems}\label{denials-and-the-pseudo-filesystems}}

Beneath the grants sits a floor of categorical denials, evaluated before
any allow, that templates carry by default and operators rarely touch.
The pseudo-filesystems are constructed rather than inherited. The kennel
gets a private tmpfs at \texttt{/tmp}, so it neither sees nor litters
the host's; its \texttt{/proc} is the fresh one of its own PID
namespace, showing only its own processes, with the host's hidden rather
than filtered; and \texttt{/dev} is assembled from an allowlist, with
the dangerous nodes denied by default, the memory devices, the GPU,
audio, video, and input devices among them, present only where a policy
passes one through deliberately. The kennel is given the device and
pseudo-filesystem surface it needs to run and nothing of the host's that
it does not, which is the same absence the home is built on, applied to
the rest of the tree.

\hypertarget{the-trust-manifest}{%
\section{6.5 The trust manifest}\label{the-trust-manifest}}

One filesystem threat is not about what the kennel reads but about what
it writes, and it reaches past the kennel's own life. A workload with a
writable bind into a project can plant a deferred trigger, a build
recipe, a repository hook, an editor task, that does nothing inside the
kennel and fires later in the operator's own unconfined shell (T2.8).
The defence is a host-readable marker the confined workload cannot
forge. A \texttt{.trust-manifest.json} at the root of each writable
workspace pins the hash of every known trigger:

\begin{kennelcode}
{
  "version": "2.0",
  "generator": "kennel compile",
  "execution": {
    "triggers": {
      "Makefile":              { "kind": "content", "sha256": "sha256:9f2b1c..." },
      ".git/hooks/pre-commit": { "kind": "content", "sha256": "sha256:4c81a0..." }
    }
  }
}
\end{kennelcode}

Each trigger path maps to the pinned hash of its bytes, which addresses
a copy held in a store beside the manifest; the record carries more than
the hash, the file mode so a revert cannot silently drop a setuid bit
and the provenance of the pin, but the hash is what a divergence is
caught against. The manifest is written by the CLI before the daemon is
contacted, and refreshed only by the deliberate \texttt{kennel\ review},
never as a side effect of a run that would otherwise launder the
workload's own edits. The parser for it runs in the CLI; the daemon
never links it.

Keeping the workload from forging the pins is a construction problem,
and it is solved with the same over-mount the view is built from. Inside
every writable bind, after the binds and before the pivot, the spawn
over-mounts an empty read-only file at the manifest's path. Inside the
kennel the manifest is present but empty: the workload can neither read
the pins to learn what is watched nor write them to cover a change. The
host inode underneath is untouched, so the operator's tools read the
real manifest. This is deliberately not a Landlock denial, because
Landlock cannot make a file in a granted directory read as empty while
its parent stays readable; the view over-mount is the only mechanism
that delivers the no-read property, and it is applied here to one
well-known path. The pinned bytes themselves are kept beside the
manifest in a content-addressed, operator-owned store, masked the same
way, which is what lets review show a diff and restore a baseline rather
than only detect a change.

The manifest is enforced from two sides. While the kennel runs, the
daemon watches the pinned triggers with an unprivileged \texttt{inotify}
in the operator's context and acts on a mutation through the cgroup it
already holds, warning, freezing the workload, or killing it, the one
place the filesystem defence reacts on the workload live rather than
only informing the host. After the run, cooperating host tooling refuses
to execute a trigger whose hash has diverged from its pin, and
\texttt{kennel\ review} shows the operator a unified diff and re-pins on
approval, the human sign-off the masked workload cannot perform; a
revert restores each changed trigger from its stored bytes and removes a
planted one. The masking and the pins are structurally complete, the
workload provably cannot read or forge them, and the honest residual is
the catalogue: the defence is only as wide as the set of trigger kinds
it knows to watch. For the writable bind that needs no concurrent host
access, marking it exclusive shadows the source host-side for the run,
so the operator and the workload cannot share the path while the kennel
lives, closing the live channel the enumerated catalogue cannot.
